\subsubsection*{Incomplete Connective Sets}

% BEGIN GENERATED ARTIFACT: def:incomplete-connective-set-propositional-logic
\begin{definitionbox}{Definition (Incomplete Connective Set)}
\begin{definition}[Incomplete Connective Set]
\label{def:incomplete-connective-set-propositional-logic}
A connective basis is \emph{incomplete} if it is not functionally complete.
\end{definition}
\end{definitionbox}

\begin{remark*}[Standard quantified statement]
\[
\operatorname{Incomplete}(B)
\Longleftrightarrow
\neg\operatorname{FunctionallyComplete}(B).
\]
\end{remark*}

\begin{remark*}[Definition predicate reading]
\(\operatorname{Incomplete}(B)\) means that some Boolean function cannot be
represented by any formula using only connectives from \(B\).
\end{remark*}

\begin{remark*}[Interpretation]
Incompleteness is usually proved by finding an invariant preserved by every
formula built from the basis. If some Boolean function fails that invariant, the
basis cannot express every Boolean function.
\end{remark*}

\begin{dependencies}
\begin{itemize}
  \item \hyperref[def:functionally-complete-connective-set-propositional-logic]{Functionally Complete Connective Set}
  \item \hyperref[def:boolean-function-propositional-logic]{Boolean Function}
\end{itemize}
\end{dependencies}
% END GENERATED ARTIFACT: def:incomplete-connective-set-propositional-logic

% BEGIN GENERATED ARTIFACT: def:monotone-boolean-function-propositional-logic
\begin{definitionbox}{Definition (Monotone Boolean Function)}
\begin{definition}[Monotone Boolean Function]
\label{def:monotone-boolean-function-propositional-logic}
A Boolean function \(f:\mathbb{B}^n\to\mathbb{B}\) is \emph{monotone} if
changing input coordinates from \(\False\) to \(\True\) never changes the output
from \(\True\) to \(\False\).
\end{definition}
\end{definitionbox}

\begin{remark*}[Standard quantified statement]
\[
\operatorname{Monotone}(f)
\Longleftrightarrow
(\forall a,b\in\mathbb{B}^n)
\bigl(a\le b\to f(a)\le f(b)\bigr),
\]
where \(\False\le\True\) coordinatewise.
\end{remark*}

\begin{remark*}[Definition predicate reading]
\(\operatorname{Monotone}(f)\) means that increasing truth inputs cannot make the
output decrease from true to false.
\end{remark*}

\begin{remark*}[Interpretation]
The connectives \(\land\) and \(\lor\) are monotone. Negation is not monotone,
because changing \(P\) from false to true changes \(\neg P\) from true to false.
\end{remark*}

\begin{dependencies}
\begin{itemize}
  \item \hyperref[def:boolean-function-propositional-logic]{Boolean Function}
\end{itemize}
\end{dependencies}
% END GENERATED ARTIFACT: def:monotone-boolean-function-propositional-logic

% BEGIN GENERATED ARTIFACT: thm:conjunction-disjunction-not-functionally-complete-propositional-logic
\begin{theorembox}{Theorem (Conjunction and Disjunction are not Functionally Complete)}
\begin{theorem}[Conjunction and Disjunction are not Functionally Complete]
\label{thm:conjunction-disjunction-not-functionally-complete-propositional-logic}
The connective basis \(\{\land,\lor\}\) is not functionally complete.

\noindent\hyperref[prf:conjunction-disjunction-not-functionally-complete-propositional-logic]{Go to proof.}
\end{theorem}
\end{theorembox}

\begin{remark*}[Standard quantified statement]
\[
\neg\operatorname{FunctionallyComplete}(\{\land,\lor\}).
\]
\end{remark*}

\begin{remark*}[Predicate reading]
There is at least one Boolean function that cannot be represented by any formula
using only conjunction and disjunction.
\end{remark*}

\begin{remark*}[Interpretation]
Every formula built using only \(\land\) and \(\lor\) represents a monotone
Boolean function. Since negation is not monotone, no formula using only
\(\land\) and \(\lor\) can represent negation. Hence the basis is incomplete.
\end{remark*}

\begin{dependencies}
\begin{itemize}
  \item \hyperref[def:incomplete-connective-set-propositional-logic]{Incomplete Connective Set}
  \item \hyperref[def:monotone-boolean-function-propositional-logic]{Monotone Boolean Function}
  \item \hyperref[def:functionally-complete-connective-set-propositional-logic]{Functionally Complete Connective Set}
\end{itemize}
\end{dependencies}
% END GENERATED ARTIFACT: thm:conjunction-disjunction-not-functionally-complete-propositional-logic
